\documentclass[11pt]{article}
\usepackage{amsfonts}
\usepackage[T1]{fontenc}
\usepackage{mathabx,graphicx}


\newcommand{\test}{\circlearrowright}
\def \loop {\ensuremath{\rotatebox[origin=c]{-90}{$\circlearrowright$}}}
\def \nestedloop {\ensuremath{\rotatebox[origin=c]{-90}{$\circlearrowright$}}^n}

\begin{document}


\section*{Ch. 3 Functions}

\section{Definition of Output}
Define output ->, "Return"

\section{Definition of Function}
Define function $\Phi$ as a map from "input", set $\{S_{in}\}$ to "output" set $\{S_{out}\}$
\begin{center}
$a_i \in  S_1,  i=1,2,...,n-1,n$\\
$b_i \in  S_2,  i=1,2,...,m-1,m$\\
$\Phi[n] \equiv $
\end{center}

\subsection{Definition of Complete Function}
\begin{center}
$a_i \in  S_1,  i=1,2,...,n-1,n$\\
$b_i \in  S_2,  i=1,2,...,m-1,m$\\
\end{center}


\subsection{Definition of Incomplete Function}
\begin{center}
$a_i \in  S_1,  i=1,2,...,n-1,n$\\
$b_i \in  S_2,  i=1,2,...,m-1,m$\\
\end{center}






\section{Definition of Mux "Multiplexor"}
Define a mux "Multiplexor" as a map from "inputs" elements $\{a_{i}$,...\} i>0 to "output" sets $\{S_j,...\}$ j > 0

\subsection{Definition of Complete Mux}
\begin{center}
$a_i \in  S_1,  i=1,2,...,n-1,n$\\
$b_i \in  S_2,  i=1,2,...,m-1,m$\\
\end{center}


\subsection{Definition of Incomplete Mux}
\begin{center}
$a_i \in  S_1,  i=1,2,...,n-1,n$\\
$b_i \in  S_2,  i=1,2,...,m-1,m$\\
\end{center}


\subsection{A function is always a mux, a mux is not always a function}







\section{Definition of a System}
Define a System as a tensor of functions, capable of mapping "inputs" set $\{S_{i}$,...\} i>0 to "outputs" set $\{S_j,...\}$ j>0


\subsection{Definition of Complete System}
\begin{center}
$a_i \in  S_1,  i=1,2,...,n-1,n$\\
$b_i \in  S_2,  i=1,2,...,m-1,m$\\
\end{center}


\subsection{Definition of Incomplete System}
\begin{center}
$a_i \in  S_1,  i=1,2,...,n-1,n$\\
$b_i \in  S_2,  i=1,2,...,m-1,m$\\
\end{center}



\section{Definition of Dimensionality}
Research Vector\\
Research Vector Space, other definitions of dimensionality



\section{Citations}
$\lbrack 1 \rbrack \hspace{2mm}$ https://en.wikipedia.org/wiki/Set\_(mathematics)\\
$\lbrack 2 \rbrack \hspace{2mm}$ https://en.wikipedia.org/wiki/Complement\_(set\_theory)\#Relative\_complement



\end{document}